The introduced \CPS model allows us to capture appearance, geometry and shape information of parts and pairs of parts in the final level of the cascade---much richer than the standard geometric deformation costs and texture filters of previous PS models~\cite{felz05,devacrf,ferrari08,andriluka09}.  
%Table~\ref{feat_table} lists all features that we use and will describe in this section.  
Each part is modeled as a rectangle anchored at the part joint with the major axis defined as the line segment between the joints (see Figure~\ref{fig:ps}).  For training and evaluation, our datasets have been annotated only with this part axis.

%The side of a part is defined as the longer side of part rectangle.  
%Further, the major axis of a part is the line segment starting at the part joint parallel to the part side.

%In the following, we will use the support of each part $l_i$ in the image defined as a rectangle anchored at the part joint position position $(l_{ix}, l_{iy})$ and aligned with the part orientation $(l_{iu}, l_{iv})$  (see Fig.~\ref{fig:ps}). The dimensions $(h_i, w_i)$ of the support rectangle for each part are predefined, where the width is set $0.25$ of the length. The side of the part is defined as the longer side of the supporting rectangle. Further, the main axis of a part is the line segment starting at the part joint parallel to the part side.

\begin{figure}[t!]
\begin{center}
\includegraphics[width=0.9\textwidth]{figs/empty.jpg}
\caption{\label{cc_fig} Left: input image; Middle left: segmentation with segment boundaries and their touching points in red. 
%A long sequence of boundaries form a contour if the angle at each touching point is small.
Middle right: contour edges which support part $y_i$ and have normals which do not deviate from the part axis normal by more than $\omega$. 
%Small $\omega$ results in edges which are parallel and close to the part axis.
Right: first and second order moments of the region lying under the major part axis.}
\end{center}
\vskip -0.4in
\end{figure}
\mypar{Shape}
We express the shape of limbs via region and contour information. We use contour cues to capture the notion that limbs have a long smooth outline connecting and supporting both the upper and lower parts.  Region information is used to express coarse global shape properties of each limb, attempting to express the fact the limbs are often supported by a roughly rectangular collection of regions---the same notion that drives the bottom-up hypothesis generation in~\cite{mori04,Srinivasan07}.
% of segments A second shape feature captures the notion that arms tend to have the shape of rectangles.

\mypar{Shape/Contour} We detect long smooth contours from sequences of image segmentation boundaries obtained via NCut~\cite{cour05}.
%sequences of boundaries of image segments.
We define a graph whose nodes are all boundaries between segments with edges linking touching boundaries. Each contour is a path in this graph (see Fig.~\ref{cc_fig}, middle left). To reduce the number of possible paths, we restrict ourselves to all shortest paths. To quantify the smoothness of a contour, we compute an angle between each two touching segment boundaries\footnote{This angle is computed as the angle between the lines fitted to the segment boundary ends, defined as one third of the boundary. }. The smoothness of a contour is quantified as the maximum angle between boundaries along this contour. Finally, we find among all shortest paths those whose length exceeds $\ell_{\textrm{th}}$ pixels and whose smoothness is less then $s_{\textrm{th}}$ and denote them by $\{c_1,\dots c_m\}$.\footnote{We set $\ell_{\textrm{th}} = 60$ pixels, $s_{\textrm{th}} = 45^\circ$ resulting in $15$ to $30$ contours per image.}

% To define the above features, we extract segments using NCut \cite{cour05}. For the contour features, we extract long smooth contours, each defined as a sequence of linked segment boundaries, such that the transitions between boundaries are smooth (see Fig.~\ref{cc_fig}, middle left). To obtain such sequences, we define a graph, whose nodes are all boundaries between segments with edges linking touching boundaries. For each pair of segment boundaries, we compute the shortest path between them, which represents a contour in the image. We retain contours which are longer than $l_{\textrm{th}}$ pixels and are smooth. To quantify the smoothness of a contour, we compute an angle between each two touching segment boundaries. This can be done by fitting a line to the segment boundary end, defined as one third of the boundary, and compute the angle between the the boundary ends at the touching point.  The smoothness of a contour is quantified as the maximum angle between boundaries along this contour. We retain contours whose smoothness is less than $s_{\textrm{th}}$. In our implementation we use $30$ segments and set $l_{\textrm{th}} = 60$ pixels, $s_{\textrm{th}} = 45^\circ$, which results in $m$ contours per image $\{c_1,\dots c_m\}$ for $m$ usually between $15$ and $30$.

We can use the above contours to define features for each pair of lower and upper arms, which encode the notion that those two parts should share a long smooth contour, which is parallel and close to the part boundaries. For each arm part $l_i$ and a contour $c_k$ we can estimate the edges of $c_k$ which lie inside one of the halves of the supporting rectangle of $l_i$ and whose edge normals build an angle smaller than $\omega$ with the normal of the part axis (see Fig.~\ref{cc_fig}, right). We denote the number of those edges by $q_{ik}(\omega)$. Intuitively, a contour supports a limb if it is mostly parallel and enclosed in one of the limb sides, i.e.~the value $q_{ik}(\omega)$ is large for small angles $\omega$. A pair of arm limbs $l_i$, $l_j$ should have a high score if both parts are supported by a contour $c_k$, which can be expressed as the following two scores
\[
  \textrm{cc}_{ijk}^{(1)}(\omega, \omega') = \frac{1}{2}\left(\frac{q_{ik}(\omega)}{h_i} + \frac{q_{jk}(\omega')}{h_j}\right)\quad\textrm{and}\quad\textrm{cc}_{ijk}^{(2)}(\omega, \omega') = \min\left\{\frac{q_{ik}(\omega)}{h_i}, \frac{q_{jk}(\omega')}{h_j}\right\}
\]
where we normalize $q_{ik}$ by the length of the limb $h_i$ to ensure that the score is in $[0,1]$. The first score measures the overall support of the parts, while the second measures the minimum support. Hence, for $l_i$, $l_j$ we can find the highest score among all contours, which expresses the highest degree of support which this pair of arms can receive from any of the image contours:
\[
\textrm{cc}_{ij}^{(t)}(\omega, \omega') = \max_{k\in\{1, \dots , m\}}\textrm{cc}_{ijk}^{(t)}(\omega, \omega'), \quad\textrm{for} \quad t\in\{1,2\}
\]
By varying the angles $\omega$ and $\omega'$ in a set of admissible angles $\Omega$ defining parallelism between the part and the contour, we obtain $|\Omega|^2$ contour features\footnote{We set $\Omega=\{10^\circ, 20^\circ, 30^\circ\}$, which results in $18$ features for both scores.}.
% For the left limbs we use the right half of the supporting rectangle, viewed from the part joint, while for the right limbs we use the left half. This corresponds to the outer contour of each arm

\mypar{Shape/Region Moments} We compute the first and second order moments of 
the segments lying under the major part axis (see Fig.~\ref{cc_fig}, 
right)\footnote{We select segments which cover at least $25\%$ of the part 
axis.} to coarsely express shape of limb hypotheses as a collection of 
segments, $R_i$. To achieve rotation and translation invariance, we compute the 
moments in the part coordinate system.  We include convexity information 
$|conv(R_{i})|/|R_{i}|$, where $conv(\cdot)$ is the convex hull of a set of 
points, and $|R_{i}|$ is the number of points in the collection of segments.  
We also include the number of points on the convex hull, and the number of part 
axis points that pass through $R_{i}$ to express continuity along the part 
axis. 
 
 Note that contour feature is a pure bottom-up cue which does not capture the 
precise arm shape but encodes figural properties of the arm -- smoothness and 
continuity of its boundary. It is a mid-level cue related to a configuration of 
two parts. The part shape, however, encodes a more precise shape but only of a 
single part.
 
\mypar{Appearance/Texture} Following the edge-based representation used in 
\cite{latentsvm}, we model the appearance the body parts using Histogram of 
Gradient (HoG) descriptor.  For each of the 6 body parts -- head, torso, upper 
and lower arms -- we learn an individual Gentleboost classifier 
\cite{friedman00} on the HoG features using the Limbs Annotated in Movies 
Dataset\footnote{LAMDa is available at 
\textit{http://vision.grasp.upenn.edu/video}}.  %The output of the classifier 
is our first appearance feature.

\mypar{Appearance/Color}  As opposed to HoG, color drastically varies between people. We use the same assumptions as \cite{eichner09} and build color models assuming a fixed location for the head and torso at run-time for each image.  We train Adaboost classifiers using these pre-defined regions of positive and negative example pixels, represented as RGB, Lab, and HSV components.  For a particular image, a 5-round Adaboost ensemble~\cite{freund1997decision} is learned for each color model (head, torso) and reapplied to all the pixels in the image.  A similar technique is also used by~\cite{strikeapose} to incorporate color.  
%Face and torso can reveal information about skin and clothing color of all other parts.  
Features are computed as the mean score of each discrimintative color model on the pixels lying in the rectangle of the part.

% on positive examples drawn from fixed rectangles, which roughly localize the person torso and head, and negative examples drawn from the background of training images. We use the color channels in RGB, Lab and HSV spaces for each pixels as features. Then we can use those models on unknown part hypotheses from the same image -- head and torso reveal information about the skin and clothing color which are same for the arm clothing and hand color. For an arm hypothesis, the two color features $f_2, f_3$ are computed as the mean score of the pixels lying in the support of the part using both models. 

We use similarity of appearance between lower and upper arms as features for 
the pairwise potentials of \CPS. Precisely, we use the $\chi^2$ distance 
between the color histograms of the pixels lying in the part support.  The 
histograms are computed using minimum-variance quantization of the RGB color 
values of each image into $8$ colors.
% Precisely, we use two features, the first $f_4$ being the $\chi^2$ distance between the color histograms of the pixels lying in the part support. The histograms are computed using minimum-variance quantization of the RGB color values of each image into $8$ colors. The second feature $f_5$ is defined in terms of the spectral embedding obtained from NCuts segmentation \cite{cour05}. NCuts is defined in terms of a pixel affinity matrix which encodes color and intervening contour similarities (see~\cite{cour05} for details). The final segmentation is obtained through discretization of the top eigenvectors (we use the top $30$) of the affinity matrix. These vectors define for each pixel an embedding vector (in our case of dimension $30$) with the property that the vectors in this embedding space are clearly clustered according to the NCut segmentation criterion and the cues used in the affinity matrix. Thus, the mean embedding vector of all pixels in the support of a body part can be interpreted as an segmentation embedding vector of the part. FInally, the dot product of two such vectors representing two arm parts can be interpreted as similarity score, which is our last appearance feature $f_5$.

\mypar{Geometry} The body part configuration is encoded in two set of features. 
The location $(l_{ix}, l_{iy})$ and orientation $(l_{iu}, l_{iv})$, included in 
the state of a part, are used added as absolute location prior features.  We 
express the relative difference between part $i$ its parent $j$ in the 
coordinate frame of the parent part as $T_{ij}(y_i) - y_j$.
% parent's coordinate frame, and $\lpart_{i\omega}$ is simply the angular 
%representation of $\lpart_i$'s direction.
%same quantities as in the original ps structure, but expanded out intThe relative positions between two parts $l_i$ and its parent $l_j$ are captured by the difference and the products of differences between the state components of the two parts: 
%\[
%f_7=(dl_x, dl_y, dl_x^2, dl_y^2, dl_xdl_y, dl_u, dl_v, dl_u^2, dl_v^2, dl_udl_v)
%\]
%where $dl_k = l_{ik} - l_{jk}$ is the difference of the state component $k$. Note that we include second order terms to model non-linear geometric interactions. 
Note we could introduce second-order terms to model a quadratic deformation cost akin to the classical PS, but we instead adopt more flexible binning or boosting of these features (see Section~\ref{implementation}).


\begin{figure*}[tb!]
\centering
\includegraphics[width=0.90\linewidth]{figs/empty.jpg}
\caption{\small \label{fig:features} Overview of features described in 
Section~\ref{sec:features}. \textbf{(a)} Two discriminative hand detector 
filters from optical flow and skin color. \textbf{(b)} Quantized color is 
matched within a frame (comparing color distributions) and over time (comparing 
$L_0$-norm patch distance). \textbf{(c)} Limb hypotheses are scored based on 
aligment to nearby contours.  \textbf{(d)} We use a few simple geometric 
features between joints in a frame, and joint persistence over time.  
\textbf{(e)} We form an estimate of foreground and background likelihood from 
dense optical flow.}
\end{figure*}

As described in Section~\ref{sec:model}, our model represents human pose and 
motion with an ensemble of tree models which capture relationships between 
different joint locations within a frame and the relationships between the same 
joint across time.  Due to the tractable nature of our tree decomposition, and 
the sparse set of states provided by a single frame cascaded 
PS~\cite{sapp10cascades}, we can afford to combine a variety of effective 
features for all unary joint and pairwise joint-joint relationships we wish to 
model.  Figure~\ref{fig:features} illustrates most of the features described 
here.


\subsection{Single-frame features}

\mypar{Geometry.} Unlike previous single-frame geometry features used in PS 
representations, we purposefully only include coarse geometric relationships 
into our model, and rely more heavily on image-based cues to estimate pose.  
This is inspired by the observation that state-of-the-art PS implementations 
learned on existing datasets tend to learn very rigid geometric priors upon 
which they rely heavily~\cite{tran10}. These models generalize poorly, 
especially to important application domains with a high degree of pose 
variation, such as action recognition.

In light of this, we use the following pairwise geometric features: (i) length 
of arms, (ii) unsigned difference in angle that the upper arm makes with 
respect to the vertical axis, (iii) difference in x-coordinate between 
left-right symmetric joints, from which our model can learn a type of repulsion 
behavior (e.g., the left shoulder should be far from the right shoulder) as 
well as a left-right order of parts (e.g., the left shoulder should be to the 
left of the right shoulder).  These features are coarsely discretized into 10 
bins.  Note that we do {\em not} express features describing the angle of the 
lower arm, leaving it free to rotate and stretch. See 
Figure~\ref{fig:features}d. 

\mypar{Color-based hand detector.}  Detecting the hand location is an extremely 
useful cue in anchoring all joint locations.  Unfortunately, traditional 
template-based part detectors such as Histogram of Gradient (HoG) detectors fail at this task due to 
hands' high variability in appearance, pose, and motion blur.  We instead learn 
a linear SVM filter on skin detection response maps computed from the publicly 
available code from~\cite{sapp10cascades} on the training data.  This can be 
evaluated efficiently using convolution at test time.  See 
Figure~\ref{fig:features}a.

\out{
 We train this filter discriminatively via linear SVM using our groundtruth 
hand locations, and it learns to encourage likely skin color in the center of 
the hand area, and discourage skin color nearby (to discourage it from firing 
in the middle of the arm, face, etc.).  See Figure~\ref{fig:features}a, bottom 
row.  Due to the linear representation, we can evaluate this detector 
efficiently at all angles via convolution.
}

\mypar{Contour support.}  For each arm joint-pair, we measure 
its support from long contours extracted in the image as follows: we take the 
number of contour points that are roughly parallel to the joint-pair line 
segment (angle less than 12 degrees) and within a spatial support region 
(approximately 25\% of the length of the average groundtruth limb).  We then 
use as a feature the number of supporting contour points, normalized by the 
length of the limb hypothesis.  Due to the sparse contour set, 
this feature can be computed extremely efficiently, by quickly discarding 
hypotheses whose endpoints are not near any contour. See Figure~\ref{fig:features}c.

\out{
%: even when the number of limb hypotheses is large (in practice, about 150K 
%possibilities for all joint pairs), we can compute all limb contour support 
%scores in less than half a second by quickly discarding ones whose endpoints 
%are not near any contour.  
The sparse contour set also makes this a low-recall, high-precision feature, 
and complements the less localized-features well when support is found.  
}

\mypar{Color consistency.} To capture the fact that the color of pairs of 
joints is often similar due to clothing and/or skin color, we describe the 
color consistency of pairs of joints via the $\chi^2$-distance between color 
histograms obtained from small image patches (with side length 10\% of image 
dimensions) extracted around each joint.  See Figure~\ref{fig:features}b.

\mypar{Figure from flow.}
We use several features based on dense optical flow from~\cite{optflow}, which we compute between 
adjacent frames. We obtain a rough estimate of the foreground of each clip by 
assuming there are only 2 planes of motion: foreground (figure) and background.  
Given a detected person, we estimate the background motion by computing the 
median flow vector $\mu_{bg}$ outside the detected person bounding box, and not 
considering outlier flow with magnitude greater than the $75^{th}$ percentile.  
We then subtract off $\mu_{bg}$ from the flow field and take the magnitude of 
the flow as an estimate of foreground likelihood, as in 
Figure~\ref{fig:features}e.  We incorporate this as a unary feature for each 
joint, and a pairwise feature by computing the average sampled evenly along 
the line segment between joint pairs.

\mypar{Flow-based hand detector.}  We exploit the fact that hands are often in 
motion (and naturally the fastest moving body part) by building a hand detector 
based on hand-shaped motion discontinuity. We extract motion discontinuities by 
computing the gradient magnitude of the flow field, and learn a linear filter 
via SVM using this motion discontinuity magnitude cue specific to hands, 
similar to the single frame hand detector based on skin-color likelihood maps.  
See Figure~\ref{fig:features}a, top row.


\mypar{Joint contour and flow support.}  We include an additional contour 
feature restricting our single-frame contour feature to only 
count support from contour points that are consistent with a large magnitude 
optical flow discontinuity.  This serves to reduce background clutter and 
restricts support to only contours that are salient both spatially and 
temporally. 

%------------------------- eccv stuff --------------------------------------
\mypar{Single-frame features from~\cite{sapp10cascades}.} We make use of 
several feature computations provided by the public implementation 
of Cascaded PS~\cite{sapp10cascades}, a state-of-the-art single frame pose estimation model; see paper for details.

We incorporate HoG limb detectors as unary features for each shoulder and elbow 
joint by taking the max response of the detector over all possible angles.  We 
also incorporate the detectors as pairwise terms for (shoulder, elbow) and 
(elbow, wrist) pairs for which we can index the detector at the appropriate 
angle.
The image-adaptive discriminative color models of clothing and skin color are 
incorporated as unary features.  Finally, we discretize our state space 
relative to the initial detected upper body into a 5x5 grid, and use membership 
in each grid cell as a feature.  This is a particularly effective feature for 
the shoulders, whose detection-relative location is relatively peaked, but not 
very informative for elbows and wrist locations in our highly articulated dataset (summarized in Figure~\ref{fig:dataset}).

\begin{figure}[]
\centering
\includegraphics[width=0.99\linewidth]{figs/empty.jpg}
\caption{\small \label{fig:dd} Sample predictions on the VideoPose2.0 test set. 
Dashed line indicates centering of torso. Magenta:
  Single-frame \cite{sapp10cascades}. Red: Dual Decomposition. Cyan:
  Single-frame agreement. Yellow: Single-variable agreement.}
\end{figure}


\subsection{Between-frame features}
\mypar{$L_0$-norm quantized color tracking.}  We capture the persistence of 
appearance over time with a simple and effective patch-based color tracker for 
each joint.  We jointly quantize each pair of consecutive images into a small 
number of color indices, using minimum variance quantization\footnote{In 
practice we use MATLAB's \texttt{rgb2ind($\cdot$)} using 32 colors}.  We then 
compare patches (of side length 10\% of the image dimensions) around the joint 
in each frame using an $L_0$-norm (Hamming loss) distance function.  Let 
$P_{t}$ and $P_{t+1}$ represent quantized color image patches around the joint 
in frame $t$ and $t+1$, with $N$ pixels.  Then the color tracking distance 
feature is $ ||P_t~-~P_{t+1}||_0~=~\frac{1}{N}~\sum_{(r,c)}~\Ind{P_{t}(r,c)~=~P_{t+1}(r,c)}$ 
where (r,c) indexes rows and columns in the patch.  This patch distance is robust to pixel value fluctuations and outliers, and encourages the {\em pixel structure} to be similar, unlike in color histogram distance tracking methods (e.g.,~\cite{ren07}).

\out{
The coarse quantization gives us a robust way to compare the color in 
consecutive frames that is tolerant to changes in lighting and pixel 
fluctuation.  The $L_0$ distance is more robust than Euclidean distance/radial 
basis similarity scores often used (e.g.,~\cite{gould08ijcv}), in which the distance grows with the distance in the color space and can thus be corrupted by a few large outlying pixels.  
Furthermore, it is a much more discriminative cue than color histogram 
distances across time (as used in, e.g.,~\cite{ren07}) because it requires that 
the {\em pixel structure} of the patch be similar in consecutive frames.  
}
%This is also in contrast to the color histogram distance we employ in a single 
%frame to compare color similarity of joints.  See Figure~\ref{fig:features}b.


\mypar{Geometry.} To express joint motion continuity, we use one simple 
geometric feature: the Euclidean distance between the joint in consecutive 
frames, discretized into 10 binary values (Figure~\ref{fig:features}d).


